\documentclass[11pt, a4paper]{article}

\input{bfw-style.tex}

\title{\vspace*{-1.5cm}\Huge\bfseries\color{bfwPrimaryDark}Aufgabenheft Session~2\\[0.4em]
  \large\color{bfwInkSoft}\textnormal{Rechte, Pflichten \& Arbeitsrecht --- Übungen im IHK-Klausurformat}}
\author{\textsf{IT Lernfeld 1 --- Das Unternehmen und die eigene Rolle im Betrieb beschreiben}}
\date{Selbstlernmaterial}

\begin{document}
\maketitle
\thispagestyle{empty}

\begin{merksatz}[title={\bfseries So nutzen Sie dieses Aufgabenheft}]
Bearbeiten Sie alle Aufgaben zunächst \emph{ohne} in die Lösungen zu schauen. Schreiben
Sie Ihre Antworten in vollständigen Sätzen --- die IHK-Prüfung verlangt formulierte
Begründungen, nicht bloße Stichworte. Beachten Sie die \emph{Operatoren} (nennen,
erläutern, beschreiben, begründen, beurteilen) --- sie geben den geforderten
Antwort-Tiefegrad vor. Erst danach öffnen Sie den Lösungsteil ab Seite~\pageref{loesungen}
und vergleichen.
\end{merksatz}

\bigskip

\textbf{Kontext:} Sie sind seit einigen Wochen als neue Auszubildende bei der \textbf{JIKU IT-Solutions GmbH},
einem IT-Dienstleistungsunternehmen mit 35 Mitarbeitenden (fiktives Übungsszenario; die Zahlen
weichen bewusst vom Lehrbuch-Standort Hamburg mit 52 Mitarbeitenden ab), tätig. Ihr Ausbilder,
Herr Schulz, informiert Sie heute über Ihre Rechte und Pflichten in der Ausbildung.

\bigskip

\noindent
\begin{tabular}{|l|l|l|l|l|l|l|l|}
\hline
\textbf{Aufgabe} & 1 & 2 & 3 & 4 & 5 & 6 & \textbf{Gesamt}\\
\hline
\textbf{Punkte} & /10 & /8 & /10 & /8 & /8 & /8 & /52\\
\hline
\end{tabular}

\tableofcontents
\vspace{1em}
\hrule

\section{Aufgaben}

\subsection*{Aufgabe~1 --- Pflichten im Ausbildungsverhältnis (10 Punkte)}

\begin{enumerate}[label=(\alph*)]
  \item \textbf{Nennen} Sie vier Pflichten, die Auszubildende nach § 13 BBiG haben. \hfill (4~P.)

  \vspace{3cm}

  \item \textbf{Erläutern} Sie, was unter der Geheimhaltungspflicht zu verstehen ist und warum sie für Auszubildende in einem IT-Unternehmen besonders wichtig ist. \hfill (3~P.)

  \vspace{3cm}

  \item \textbf{Beschreiben} Sie, was der Ausbildende nach § 14 BBiG in Bezug auf die Ausbildungsmittel verpflichtet ist zu tun. \hfill (3~P.)

  \vspace{3cm}
\end{enumerate}

\subsection*{Aufgabe~2 --- Ausbildungsvertrag und Vergütung (8 Punkte)}

\begin{enumerate}[label=(\alph*)]
  \item \textbf{Nennen} Sie vier der zehn Pflichtinhalte, die § 11 BBiG für den Berufsausbildungsvertrag vorschreibt. \hfill (4~P.)

  \vspace{3cm}

  \item Was gilt nach dem BBiG für Vereinbarungen im Ausbildungsvertrag, die dem Sinn und Zweck der Berufsausbildung widersprechen? \textbf{Begründen} Sie Ihre Antwort. \hfill (2~P.)

  \vspace{2.5cm}

  \item Wie lange wird die Ausbildungsvergütung bei unverschuldeter Erkrankung weitergezahlt und was passiert danach? \hfill (2~P.)

  \vspace{2.5cm}
\end{enumerate}

\subsection*{Aufgabe~3 --- Probezeit und Kündigung (10 Punkte)}

Herr Schulz erklärt Ihnen, dass Ihre Ausbildungskollegin Mia (17 Jahre) nach zwei Monaten Ausbildung überlegt, ob sie kündigen könnte.

\begin{enumerate}[label=(\alph*)]
  \item Wie lang darf die Probezeit im Berufsausbildungsverhältnis nach § 20 BBiG mindestens und höchstens dauern? \hfill (2~P.)

  \vspace{2cm}

  \item Mia befindet sich noch in der Probezeit. \textbf{Erläutern} Sie, welche Möglichkeiten zur Kündigung Mia und JIKU IT-Solutions in dieser Phase haben. \hfill (3~P.)

  \vspace{3cm}

  \item Angenommen, die Probezeit ist vorbei und Mia möchte nun die Ausbildung aufgeben und in einen anderen Beruf wechseln. \textbf{Beschreiben} Sie, wie sie kündigen darf und was sie dabei beachten muss. \hfill (3~P.)

  \vspace{3cm}

  \item \textbf{Begründen} Sie, ob JIKU IT-Solutions nach der Probezeit das Ausbildungsverhältnis ordentlich kündigen kann. \hfill (2~P.)

  \vspace{2.5cm}
\end{enumerate}

\subsection*{Aufgabe~4 --- Jugendarbeitsschutz (8 Punkte)}

Mia ist 17 Jahre alt. Kreuzen Sie an, ob folgende Aussagen zum JArbSchG \textbf{richtig (R)} oder \textbf{falsch (F)} sind, und korrigieren Sie falsche Aussagen in einem Satz.

\begin{center}
\begin{tabular}{p{7.5cm} c c p{4.5cm}}
\toprule
\textbf{Aussage} & \textbf{R} & \textbf{F} & \textbf{Korrektur (wenn falsch)}\\
\midrule
a) Mia darf pro Tag höchstens 8 Stunden arbeiten. & {[~]} & {[~]} & \underline{\hspace{4cm}}\\[1.5em]
b) Mia hat Anspruch auf 30 Urlaubstage pro Jahr. & {[~]} & {[~]} & \underline{\hspace{4cm}}\\[1.5em]
c) Bei einem Berufsschultag mit 6 Unterrichtsstunden ist Mia ganztägig freizustellen. & {[~]} & {[~]} & \underline{\hspace{4cm}}\\[1.5em]
d) Mehrarbeitsstunden bei Mia müssen am selben Tag ausgeglichen werden. & {[~]} & {[~]} & \underline{\hspace{4cm}}\\[1em]
\bottomrule
\end{tabular}
\end{center}
\textit{(2 Punkte je Zeile: 1 P. Kreuz korrekt, 1 P. Korrektur soweit nötig)}

\subsection*{Aufgabe~5 --- Betriebsrat und JAV (8 Punkte)}

JIKU IT-Solutions hat 35 Mitarbeitende (fiktives Übungsszenario; die Zahlen weichen bewusst vom Lehrbuch-Standort Hamburg mit 52 Mitarbeitenden ab), davon vier Auszubildende (alle zwischen 19 und 23 Jahren).

\begin{enumerate}[label=(\alph*)]
  \item \textbf{Erklären} Sie, ob bei JIKU IT-Solutions ein Betriebsrat gewählt werden kann, und begründen Sie Ihre Antwort. \hfill (3~P.)

  \vspace{3cm}

  \item \textbf{Erläutern} Sie, was die JAV ist und unter welchen Voraussetzungen sie bei JIKU IT-Solutions eingerichtet werden kann. \hfill (3~P.)

  \vspace{3cm}

  \item Nennen Sie \textbf{zwei Aufgaben} der JAV im Betrieb. \hfill (2~P.)

  \vspace{2cm}
\end{enumerate}

\subsection*{Aufgabe~6 --- Fallaufgabe: JIKU IT-Solutions (8 Punkte)}

Herr Schulz bittet Sie, eine betriebliche Situation zu beurteilen: Ein Auszubildender bei JIKU IT-Solutions wurde dauerhaft in der Poststelle eingesetzt, weil dort Personalmangel herrschte. Die Tätigkeit hat keinen Bezug zur IT-Ausbildung.

\begin{enumerate}[label=(\alph*)]
  \item \textbf{Beurteilen} Sie, ob dieses Vorgehen mit dem BBiG vereinbar ist, und nennen Sie den relevanten Paragrafen. \hfill (4~P.)

  \vspace{3.5cm}

  \item Der Auszubildende möchte wissen, an wen er sich wenden kann, um die Situation zu ändern. \textbf{Nennen} Sie zwei mögliche Ansprechpartner und \textbf{erläutern} Sie deren Aufgabe. \hfill (4~P.)

  \vspace{3.5cm}
\end{enumerate}

\vspace{1em}
\noindent\textbf{Punkte gesamt: 52}

\newpage
\section*{Lösungshinweise für Lehrende}\label{loesungen}
\addcontentsline{toc}{section}{Lösungshinweise für Lehrende}

\subsection*{Lösung Aufgabe~1}
\begin{loesung}
\textbf{(a)} Vier Pflichten nach § 13 BBiG (je 1 Punkt, beliebige vier aus):
Lernpflicht, Teilnahmepflicht (Berufsschule und Ausbildungsmaßnahmen), Weisungsgebundenheit,
Einhaltung der betrieblichen Ordnung, Führung des Ausbildungsnachweises,
pflegliche Behandlung der Ausbildungsmittel, Benachrichtigungspflicht bei Abwesenheit,
Geheimhaltungspflicht.

\textbf{(b)} Die Geheimhaltungspflicht nach § 13 BBiG verpflichtet Auszubildende, über
Betriebs- und Geschäftsgeheimnisse Stillschweigen zu bewahren. Für ein IT-Unternehmen
wie JIKU IT-Solutions ist dies besonders bedeutsam, da Auszubildende häufig Zugang zu
sensiblen Kundendaten, Systemkonfigurationen und internen Projekten erhalten. Ein Verstoß
könnte erheblichen wirtschaftlichen Schaden oder Vertrauensverlust verursachen.

\textbf{(c)} Nach § 14 BBiG sind Ausbildende verpflichtet, Auszubildenden kostenlos alle
Ausbildungsmittel zur Verfügung zu stellen, die für die Berufsausbildung und das Ablegen von
Prüfungen erforderlich sind (z.\,B.\ Hardware, Software, Fachbücher). Außerdem müssen sie
die ordnungsgemäße Führung des Ausbildungsnachweises überwachen.
\end{loesung}

\subsection*{Lösung Aufgabe~2}
\begin{loesung}
\textbf{(a)} Vier beliebige der zehn Pflichtinhalte nach § 11 BBiG (je 1 Punkt), z.\,B.:
Art, Dauer und Ziel der Ausbildung; Probezeit; Höhe und Zahlung der Vergütung; Dauer des
Urlaubs; Form des Ausbildungsnachweises; Beginn und Dauer der Ausbildung.

\textbf{(b)} Vereinbarungen, die den Vorschriften des BBiG widersprechen oder dem Sinn und
Zweck der Berufsausbildung zuwiderlaufen, sind \emph{nichtig} --- d.\,h.\ von Anfang an
unwirksam. Das gilt unabhängig davon, ob beide Seiten zugestimmt haben. Das BBiG schützt
Auszubildende damit auch vor nachteiligen Vertragsklauseln.

\textbf{(c)} Die Ausbildungsvergütung wird bei unverschuldeter Erkrankung nach dem
Entgeltfortzahlungsgesetz (EntgFG) für bis zu \textbf{sechs Wochen} weitergezahlt.
Danach zahlt die gesetzliche Krankenkasse \textbf{Krankengeld}.
\end{loesung}

\subsection*{Lösung Aufgabe~3}
\begin{loesung}
\textbf{(a)} Die Probezeit beträgt nach § 20 BBiG \textbf{mindestens einen Monat} und
\textbf{höchstens vier Monate}. Eine Probezeit von fünf oder sechs Monaten wäre unzulässig.

\textbf{(b)} Während der Probezeit kann \emph{jede Seite} ohne Einhaltung einer Frist
und ohne Angabe von Gründen kündigen (§ 22 Abs.~1 BBiG). Mia könnte also fristlos
kündigen, ebenso wie JIKU IT-Solutions.

\textbf{(c)} Nach der Probezeit kann Mia das Ausbildungsverhältnis kündigen, wenn sie
die Berufsausbildung aufgeben oder einen anderen Beruf anstreben möchte. Die Kündigung
muss \textbf{schriftlich}, mit \textbf{Angabe des Grundes} und mit einer
\textbf{Frist von vier Wochen} erfolgen (§ 22 Abs.~2 Nr.~2 BBiG). \textit{Hinweis: Anders
als der Vertrag (Textform genügt) muss die Kündigung stets schriftlich erfolgen
(§~22 Abs.~3 BBiG).}

\textbf{(d)} Nein. § 22 BBiG erlaubt Ausbildenden nach der Probezeit \emph{keine} ordentliche
Kündigung. Nur eine außerordentliche fristlose Kündigung aus \textbf{wichtigem Grund} ist
möglich --- und muss innerhalb von zwei Wochen nach Kenntnis des Grundes ausgesprochen werden.
\end{loesung}

\subsection*{Lösung Aufgabe~4}
\begin{loesung}
(a) \textbf{Richtig} --- § 8 JArbSchG begrenzt die tägliche Arbeitszeit auf höchstens 8 Stunden (in Ausnahmen 8{,}5 Stunden).\\[0.4em]
(b) \textbf{Falsch} --- Mia ist 17 Jahre alt (noch nicht 18 Jahre), daher hat sie Anspruch auf \textbf{25 Urlaubstage} (§ 19 JArbSchG: noch nicht 18 Jahre → 25 Tage; noch nicht 17 Jahre → 27 Tage; noch nicht 16 Jahre → 30 Tage).\\[0.4em]
(c) \textbf{Richtig} --- § 9 JArbSchG schreibt die ganztägige Freistellung vor, wenn der Berufsschulunterricht mehr als fünf Unterrichtsstunden umfasst (einmal pro Woche).\\[0.4em]
(d) \textbf{Falsch} --- Mehrarbeit muss innerhalb der folgenden \textbf{drei Wochen} (nicht am selben Tag) durch Freizeit ausgeglichen werden.
\end{loesung}

\subsection*{Lösung Aufgabe~5}
\begin{loesung}
\textbf{(a)} Ja, bei JIKU IT-Solutions kann ein Betriebsrat gewählt werden. Das
Betriebsverfassungsgesetz (BetrVG) ermöglicht die Wahl in Betrieben mit mindestens
\textbf{fünf wahlberechtigten Beschäftigten} (ab 18 Jahre). JIKU IT-Solutions hat 35
Mitarbeitende, erfüllt damit die Voraussetzung deutlich. Die Amtszeit des gewählten
Betriebsrats beträgt vier Jahre.

\textbf{(b)} Die JAV (Jugend- und Auszubildendenvertretung) ist die gewählte Interessenvertretung
aller Jugendlichen und Auszubildenden im Betrieb. Voraussetzung für ihre Einrichtung sind
mindestens \textbf{fünf wahlberechtigte Jugendliche oder Auszubildende}. Da JIKU IT-Solutions
vier Auszubildende zwischen 19 und 23 Jahren hat, wäre eine JAV erst bei einem weiteren
wahlberechtigten Auszubildenden oder Jugendlichen möglich.

\textbf{(c)} Zwei Aufgaben der JAV (je 1 Punkt):
(1) Überwachung der Einhaltung von BBiG, JArbSchG und Betriebsvereinbarungen;
(2) Weiterleitung von Verbesserungsvorschlägen und Beschwerden der Auszubildenden über
den Betriebsrat an den Arbeitgeber.
\end{loesung}

\subsection*{Lösung Aufgabe~6}
\begin{loesung}
\textbf{(a)} Das Vorgehen ist \emph{nicht} mit dem BBiG vereinbar. § 14 BBiG schreibt vor,
dass Ausbildende Auszubildenden ausschließlich Aufgaben übertragen dürfen, die dem
\textbf{Ausbildungszweck dienen} und ihren körperlichen Kräften angemessen sind. Der
dauerhafte Einsatz in der betriebsfremden Poststelle ohne IT-Bezug verstößt gegen § 14 BBiG,
da er dem Ausbildungsziel des Fachinformatikers nicht dient. Darüber hinaus gefährdet er
das Erreichen des Ausbildungsziels.

\textbf{(b)} Zwei mögliche Ansprechpartner:
(1) \emph{Betriebsrat}: Der Betriebsrat hat ein Mitbestimmungsrecht bei Fragen der
Berufsausbildung und kann beim Arbeitgeber auf Abhilfe drängen.
(2) \emph{IHK (zuständige Kammer)}: Als zuständige Stelle nach BBiG berät die IHK
Auszubildende bei Problemen mit dem Ausbildenden und kann vermitteln oder ein
Schlichtungsverfahren einleiten.
\end{loesung}

\end{document}
